Para \(a\ne\pm1\),
\[x_1=Ce^{at}+De^{-t},\quad x_2=aCe^{at}-De^{-t},\quad
x_3=Ee^t-\frac{a^2C}{a-1}e^{at}+\frac D2e^{-t}.\]
Para \(a=1\), las dos primeras expresiones siguen siendo válidas y
\(x_3=(E-Ct)e^t+\tfrac D2e^{-t}\).
Para \(a=-1\),
\[x_1=(C+Dt)e^{-t},\quad x_2=(D-C-Dt)e^{-t},\quad
x_3=Ee^t+\left(\frac C2-\frac{3D}{4}+\frac{Dt}{2}\right)e^{-t}.\]
Las constantes \(C,D,E\) son reales y \(t\in\mathbb R\).
